\documentclass[10pt,a4paper]{amsart}
\usepackage[margin=19mm]{geometry}
\usepackage{amsmath,amssymb,mathtools}
\usepackage[scr=rsfs,cal=euler]{mathalfa}
\usepackage{xfrac}
\usepackage[unicode,hidelinks,bookmarks=false]{hyperref}
\newcommand{\ie}{i.e.\ }
\newcommand{\cf}{cf.\ }
\newcommand{\resp}{respectively\ }
\newcommand{\Subsection}{\subsection}
\providecommand{\nobreakdash}{\nobreak\hbox{-}}
\setlength{\emergencystretch}{2em}
\multlinegap0pt
\usepackage{bm}
\usepackage{bbm}
\usepackage{dsfont}
\usepackage{xspace}
\usepackage{cancel}
\usepackage{mathtools}
\usepackage{amsthm}
\makeatletter
\@ifundefined{theorem}{\newtheorem{theorem}{Theorem}}{}
\@ifundefined{lemma}{\newtheorem{lemma}[theorem]{Lemma}}{}
\makeatother
\newcommand{\N}{{\mathbb N}}
\newcommand{\Q}{{\mathbb Q}}
\newcommand{\mc}{\mathcal}
\title[Generalized Thue-Morse measures: spectral and fractal analys]{Generalized Thue-Morse measures: spectral and fractal analysis}
\author{Philipp Gohlke, Marc Kesseb\"ohmer, Tanja I. Schindler}
\date{Source: arXiv:2509.22109v1 (CC BY 4.0)}
\begin{document}
\maketitle

\noindent\emph{Source and licence.} Generalized Thue-Morse measures: spectral and fractal analysis. Version arXiv:2509.22109v1 (CC BY 4.0), distributed under the Creative Commons Attribution 4.0 International licence, as recorded in the arXiv licence metadata for this exact version and shown on the abstract page https://arxiv.org/abs/2509.22109v1. Full rights notice: licenses/arxiv-2509.22109-CC-BY-4.0.txt.\par\smallskip

\noindent\textit{Editorial scope.} Counted unit: the Lemma whose source label is \texttt{LEM:block-for-variation-bound} in the article (source0.tex lines 2665--2672, source label \texttt{LEM:block-for-variation-bound}), delivered together with the article's own complete original proof of it (source0.tex lines 2674--2793, one \texttt{proof} environment of 120 lines). No mathematics is written, repaired or re-derived: statement and proof are the article's own text line for line. Required context retained uncounted: LEM:word-power (lines 2117--2120); LEM:hitting-depths-relations (lines 2158--2161); LEM:periodic-hitting-depth (lines 2169--2172); LEM:hitting-word-structure (lines 2190--2204); LEM:closing-extension-II (lines 2283--2291); LEM:s-hitting-relation (lines 2371--2374); LEM:psi-distance-bound (lines 2506--2517); LEM:psi-on-hitting (lines 2531--2545). Every definition, display and label of those lines is retained unchanged. Omitted from this package and not redistributed: 1--2116, 2121--2157, 2162--2168, 2173--2189, 2205--2282, 2292--2370, 2375--2505, 2518--2530, 2546--2664, 2673--2673, 2794--3566. Nothing inside a retained excerpt is omitted or altered: every retained body is delivered exactly as the article's lines are written, so no omission receipt of any kind accompanies this package. Citations: the retained excerpts carry no citation command at all, so no bibliography entry of the article is transcribed and no reference is redistributed. Reference adaptation: none was needed and none was made; every reference inside the delivered unit resolves inside it, so no unresolved reference or citation is produced. Printed slips kept literal: no printed word, symbol, hypothesis or number of the counted statement or of its proof was altered. Layout and class: the article's own class is \texttt{amsart} with packages including pgfplots, mathrsfs, a4wide, changepage, parskip, amsmath, amssymb, bbm, color, tikz, tkz-graph, tikz-cd, dsfont, amscd; the wrapper is the lane's pinned \texttt{amsart} editorial wrapper at 10pt on A4, which restates the theorem-like environments the retained text needs and adds the article's own macro definitions that the retained text needs (\texttt{\textbackslash N}, \texttt{\textbackslash Q}, \texttt{\textbackslash mc}), with extra wrapper packages (bm, bbm, dsfont, xspace, cancel, mathtools). The delivered numbering is the unit's own. Assets: no retained span contains a figure environment, a graphics-inclusion command, a PDF-page inclusion, a table or a TikZ picture; no image file is copied into this package, the package declares no asset dependency and no image file is redistributed with it. Licence: version arXiv:2509.22109v1 of this article is distributed under the Creative Commons Attribution 4.0 International licence, as recorded in the arXiv licence metadata for this exact version and shown on the abstract page https://arxiv.org/abs/2509.22109v1. A full rights notice is delivered at licenses/arxiv-2509.22109-CC-BY-4.0.txt. Characters: every retained excerpt is pure ASCII, so the delivered engine, XeLaTeX with its default Latin Modern text font, reports no missing character.
\begin{lemma}
\label{LEM:word-power}
Assume that $w = u w_p$, where $w_p$ is a prefix of $w$. Then, $w$ is a (rational) power of the word $u$.
\end{lemma}

\begin{lemma}
\label{LEM:hitting-depths-relations}
Let $w \in \Sigma^n$ and $r<n-1$ be such that $w_{[r+1,r+2]} \in \{01,10 \}$. If $k \in \widetilde{\mc H}_j(w)$ for some $j < r $, then the hitting depth of $k$ is at most $r+1$.
\end{lemma}

\begin{lemma}
\label{LEM:periodic-hitting-depth}
Assume that $\breve{c} \neq 0$ has a periodic expansion of the form $\breve{c} = u^\infty$ with $|u| = p > 1$. Then, if $w \in \Sigma^n$ and $k \in \widetilde{\mc H}_j(w)$ for some $j < n-p$, the hitting depth of $k$ is at most $j+p$.
\end{lemma}

\begin{lemma}
\label{LEM:hitting-word-structure}
Let $w \in \Sigma^n$ and let $\mc W(w) = \left\{ w_{[k,n]} : k \in \widetilde{\mc H}(w) \right\}$ be the set of full hitting words. Then there are $s = s(w) \in \N_0$, pairwise disjoint words $u_1, \ldots, u_s$, and multiplicities 
$N_1, \ldots, N_s \in \N$ with the following properties. Let $v_i = u_i^{N_i}$ for each $1\leqslant i \leqslant s$. Then,

\begin{enumerate}
\item\label{en: 1} $\mc W(w) = \bigcup_{i=1}^s \mc W_i$, where
\[
\mc W_i = \left\{ u_i^j v_{i-1} \cdots v_1: 1 \leqslant j \leqslant N_{i} \right\}.
\]
\item\label{en: 2} Each element of $\mc W_i$ is a rational power of $u_i$.
\item\label{en: 3} $|u_{i+1}| \geqslant \max_{w \in \mc W_i}|w|/2$ for all $1\leqslant i < s$.
\end{enumerate}
In particular, $\# \widetilde{\mc H}(w) = \sum_{i=1}^s N_i$ and $s \leqslant \lceil \log_{3/2} n \rceil$.
\end{lemma}

\begin{lemma}
\label{LEM:closing-extension-II}
There is a constant $C>0$ with the following property. Let $w \in \Sigma^n$ and $s = s(w)$ the number of independent hitting times in $\widetilde{\mc H}(w)$. Then there is a word $v$ such that
\begin{itemize}
\item $w_{[k,n]}v \notin \mc G$ for all $1\leqslant k \leqslant n$
\item $|v| \leqslant \max \left\{C, 2 \log_2 s \right\}$.
\end{itemize}
In particular, there is $n_0 \in \N$ such that for all $n \geqslant n_0$, we have $|v| \leqslant 2 \log_2 \log_{3/2} n$.
\end{lemma}

\begin{lemma}
\label{LEM:s-hitting-relation}
Let $w \in \Sigma^n$ and let $\delta \in (0,1)$ be such that $\delta n = \# \widetilde{\mc H}(w) - s$, with $s$ as in Lemma~\ref{LEM:hitting-word-structure}. Then we have $s \leqslant 2 + \log_{3/2} (2\delta^{-1})$.
\end{lemma}

\begin{lemma}
\label{LEM:psi-distance-bound}
For every $x$, we have
\[
2 \log (2 \rho(x,\breve{c})) \leqslant \psi^c(x)
\leqslant 2 \log (\pi \rho(x, \breve{c})).
\]
Furthermore, if $x,y \in \mathbb{T} \setminus B_r(\breve{c})$ for some $r>0$, we have the Lipschitz estimate
\[
|\psi^c(x) - \psi^c(y)| \leqslant \frac{2}{r} \rho(x,y).
\]
\end{lemma}

\begin{lemma}
\label{LEM:psi-on-hitting}
Let $w \in \Sigma^n$ and $k \in \mc H_j(w)$ for some $j < n$. If $x,y \in \langle w_{[k,n]} \rangle$, we have
\[
(k-j-1)\, 2 \log 2
\leqslant \psi^c(x) \leqslant 
(k-j +2)\, 2 \log 2
\]
and
\[
|\psi^c(x) - \psi^c(y)|
\leqslant  2^{j-n + 2}.
\]
The lower bound on $\psi^c(x)$ also holds for $j=n$.
\end{lemma}

\begin{lemma}
\label{LEM:block-for-variation-bound}
Assume the binary expansion of $\breve{c}$ to be non-periodic and let $\varepsilon > 0$. Then, there are infinitely many $n\in \N$ such that for each $w \in \mc G_n$, there is a word $v$ with $|v| \leqslant 2 \log_2 \log_{3/2} n$ and $w_{[k,n]}v \notin \mc G$ for all $1\leq k \leq n$, and such that, given $x \in \langle w \rangle$ and $y \in \langle wv \rangle$,
\[
\psi^c_n(x) - \psi^c_n(y) \leqslant \varepsilon n.
\]
The same holds for periodic $\breve{c}$, if additionally $x_{[1,n+2]} \triangleleft_p \breve{c}$.
\end{lemma}

\begin{proof}
The general strategy of proof is as follows. Given $w \in \mc G_n$, we choose $s = s(w)$ and choose $v$ as in Lemma~\ref{LEM:closing-extension-II}.
Note that
\begin{equation}
\label{EQ:psi-variation-decompose}
\psi^c_n(x) - \psi^c_n(y) = \sum_{k=1}^n \psi^c(x_k) - \psi^c(y_k)
\end{equation}
for some $x_k \in \langle w_{[k,n]} \rangle$ and $y_k \in \langle w_{[k,n]} v \rangle$. For each $k$, we can find an estimate as follows. By the choice of $v$, we have $w_{[k,n]}v \notin \mc G$, implying that $\rho(y_k,\breve{c}) \geqslant 2^{-(n + |v| + 1 - k)}$. Furthermore, we certainly have $\psi^c(x_k) \leqslant \psi^c(x_k')$ where $x_k'$ is the point in $\langle w_{[k,n]} \rangle$ maximizing the distance to $\breve{c}$, in particular $\rho(x_k', \breve{c}) \geqslant \rho(y_k,\breve{c})$. Since $\rho(x_k,y_k) \leqslant 2^{-(n+1-k)}$, we get from the Lipschitz estimate in Lemma~\ref{LEM:psi-distance-bound},
\begin{equation}
\label{EQ:psi-variation-general}
\psi^c(x_k) - \psi^c(y_k)
\leqslant \psi^c(x_k') - \psi^c(y_k)
\leqslant 2^{|v|+1}.
\end{equation} 
However, using this bound for every $k$ will provide an estimate that is too rough for our purpose.
We are therefore concerned with finding finer estimates for $\psi^c(x_k) - \psi^c(y_k)$ based on the (symbolic) hitting depth of $k$.


If $\breve{c} \neq 0$ is eventually constant, it can be written in the form $\breve{c} = u10^\infty = u01^\infty$, for some word $u$ (possibly the empty word). As before, we choose $v$ as in Lemma~\ref{LEM:closing-extension-II}, implying in particular that $w_{[k,n]}v \notin \mc G$ for all $1\leqslant k \leqslant n$. Further, we use the identity in \eqref{EQ:psi-variation-decompose} and consider different cases for $k$. Fix some $n > |u|$. Then, $\mc G_n = \{ u01^r, u10^r \}$  with $r = n-|u|-1$. If $2 \leqslant k \leqslant n - |u| - 1$, we argue that the hitting depth of $k$ is at most $k+|u|$. Indeed, in this case $w_{[k,k+|u| + 1]}$ ends with $00$ or $11$, and it can therefore not be contained in $\mc G_{|u| + 2} = \{ u01, u10\}$. 
In particular, due to Lemma~\ref{LEM:psi-on-hitting},
\[
\psi^c(x_k) - \psi^c(y_k) \leqslant 2^{|u| + 2} 2^{k-n}.
\]
Evaluating a geometric sum for such $k$, and using the general bound in \eqref{EQ:psi-variation-general} for $k=1$ and $n-|u| \leqslant k \leqslant n$, we obtain
\[
\sum_{k=1}^n \psi^c(x_k) - \psi^c(y_k)
\leqslant (|u| + 2) 2^{|v| + 1} + 4,
\]
and, due to the estimate $|v| \leqslant 2 \log_2 \log_{3/2} n$, this expression is certainly smaller than $\varepsilon n$ for large enough $n$.

Assume now that $\breve{c}$ is non-periodic and not eventually constant.
We fix four large integer constants $m, p \ll n_p \ll n'$ in a hierarchical order, with precise relations to be determined later. Since $\breve{c}$ is not eventually constant, there are infinitely many choices of words $w_1$ such that $w_1 01$ is a prefix of $\breve{c}$. We choose one such $w_1$ with $|w_1| \geqslant n'$. Let $w_2'$ be the unique word of length $m$ such that $w' = w_1 01 w_2'$ is a prefix of $\breve{c}$ and set $n = |w'|$. 
Note that every $w \in \mc G_n$ is of the form $w = w_1 01 w_2''$, $w = w_1 10 0^m$, or $w = w_1 00 1^m$. In any case, $w = w_1 u w_2$,
with $u \in \{01,10\}$ and $|w_2| \in \{m-1,m\}$, redefining $w_1 0$ as $w_1$ in the last case.
For such $w$, we choose $s=s(w)$ and $v$ as in Lemma~\ref{LEM:closing-extension-II} and use the splitting in \eqref{EQ:psi-variation-decompose}. In particular $w_{[k,n]}v \notin \mc G$ for all $1 \leqslant k \leqslant n$.
If $k \in \widetilde{H}_{j}(w)$ for some $j<|w_1|$ 
we obtain from Lemma~\ref{LEM:hitting-depths-relations} that the hitting depth $j'$ of $k$ is no more than $|w_1| + 1$,
because $w_1$ is followed by $u \in \{01,10\}$.
Hence, recalling the splitting in \eqref{EQ:psi-variation-decompose} and using that $x_{k},y_k$ are in $\langle w_{[k,n]}\rangle$, we obtain via Lemma~\ref{LEM:psi-on-hitting},
 \[
\psi^c(x_k) - \psi^c(y_k) \leqslant 2^{j'-n+2} \leqslant 2^{-m+2} \leqslant \varepsilon/4,
\]
for large enough $m$. In particular,
\[
\sum_{j < |w_1|} \sum_{k \in \widetilde{\mc H}_j(w)} \psi^c(x_k) - \psi^c(y_k) \leqslant n \varepsilon/4.
\]
Given $p$, let $n_p$ be large enough to ensure that $\breve{c}_{[1,n_p]}$ is not $q$-periodic for any $1\leqslant q \leqslant p$. 
Assume $k_i \in \widetilde{\mc H}_{j_i}(w)$, with $|w_1| \leqslant j_i < |w|$ for $i =1,2$.
If $k_1 < k_2 \leqslant |w_1|- n_p$, this requires that $w_{[k_1,|w_1|]}$ and $w_{[k_2,|w_1|]}$ are both prefixes of $\breve{c}$, and hence $w_{[k_1,|w_1|]}$ is a rational power of $w_{[k_1,k_2-1]}$, with period length $k_2 - k_1$; compare Lemma~\ref{LEM:word-power}. 
Since $w_{[k_1,|w_1|]}$ is a prefix of $\breve{c}$ of length exceeding $n_p$, this is only possible if $k_2 - k_1 > p$. Due to this minimal distance, there are at most $n/p + n_p + m  + 2$ positions $k$ of symbolic hitting depth $j > |w_1|$. Indeed, there are at most $n_p + m + 2$ choices for $|w_1| - n_p < k \leqslant n$ and for $k \leqslant |w_1| - n_p$ there are no more than $n/p$ elements due to the separation of $p$ established above.

Note that for every $\delta'>0$ we can choose first $p$ and then $n' > n_p,m$ large enough that $n>n'$ guarantees
\[
n/p + n_p + m + 2\leqslant \delta' n.
\]
For $k \in \widetilde{\mc H}_j(w)$ and $j < |w|$, we get $\psi^c(x_k) - \psi^c(y_k) \leqslant 2$, again via Lemma~\ref{LEM:psi-on-hitting}. Hence,
\[
\sum_{|w_1| \leqslant j < |w|} \sum_{k \in \widetilde{\mc H}_j(w)} \psi^c(x_k) - \psi^c(y_k) \leqslant 2 (n/p + n_p + m + 2) \leqslant n \varepsilon/4,
\]
for large enough $p$ and $n'$. 
Finally, we deal with $k \in \widetilde{\mc H}(w)$, meaning $w_{[k,n]} \triangleleft_p \breve{c}$. By the same argument as above, we obtain that $\# \widetilde{\mc H}(w) \leqslant n/p + n_p + m + 2 \leqslant \delta_0 n$, for arbitrarily small $\delta_0>0$, to be chosen later. 
Using the general bound in \eqref{EQ:psi-variation-general}, we obtain
\[
\sum_{k \in \widetilde{\mc H}(w)} \psi^c(x_k) - \psi^c(y_k) \leqslant 2^{|v| +1} \# \widetilde{\mc H}(w).
\]
Recall from Lemma~\ref{LEM:closing-extension-II} that $2^{|v|} \leqslant \max\left\{2^C, s^2\right\}$, where $s$ is the number of independent hitting times in $\widetilde{\mc H}(w)$.
If $s^2 \leqslant 2^C$, we can choose $\delta_0$ such that $\delta_0 2^C \leqslant \varepsilon /3$ and the proof if complete. If $s^2 > 2^C$, let $\delta$ be such that $\delta n = \# \widetilde{\mc H}(w) - s < \delta_0 n$. Then,
\[
2^{|v|} \# \widetilde{\mc H}(w)
\leqslant s^2 (s+ \delta n) = s^3 + s^2 \delta n.
\]
Note that again by Lemma~\ref{LEM:closing-extension-II} $s^3 \leqslant (2 \log_2 \log_{3/2} n)^3$ which is smaller than $n \varepsilon/8$ for large enough $n$. If $\delta \leqslant 0$, we bound the second term with $0$. 
If $\delta>0$, Lemma~\ref{LEM:s-hitting-relation} yields $s \leqslant 2 - \log_{3/2}(\delta/2)$ 
and therefore
\[
s^2 \delta n \leqslant  (2- \log_{3/2}(\delta/2) )^2 \delta n \leqslant n \varepsilon/8,
\]
by choosing $\delta_0$ small enough and recalling $\delta < \delta_0$. In any case,
\[
\sum_{k \in \widetilde{\mc H}(w)} \psi^c(x_k) - \psi^c(y_k) \leqslant n \varepsilon/2,
\]
and the proof is complete in case that $\breve{c}$ is not periodic.

In the last step, we consider the case that the binary representation of $\breve{c}$ is periodic, of the form $\breve{c} = u^\infty$ for some $u \in \Sigma^p$ and $p \in \N$.
Without loss of generality, the length of $u$ is minimal with this property. Consider $n = mp$ for some $m \in \N$ and note that, by assumption, $x \in \langle w v' \rangle$ with $v' \in \Sigma^2$ and $wv'$ a prefix of $\breve{c}$, hence a power of $u$. 
In particular, $w=u^m$ and $v'$ is a rational power of $u$. We prove that the claimed estimate holds for some $v$ that starts with some fixed, arbitrary word in $\Sigma^2 \setminus \mc G_2$. 
Note that 
\[
\psi^c_n(x) - \psi^c_n(y) = \sum_{k=1}^n \psi^c(x_k) - \psi^c(y_k)
\]
for some $x_k \in \langle w_{[k,n]} v' \rangle$ and $y_k \in \langle w_{[k,n]} v \rangle$.
For $k$ of the form $k = rp+1$ with $0\leqslant r < m$, we have that $w_{[k,n]} v' = u^q$ (for some $q \in \Q$) is a prefix of $\breve{c}$. Hence, for such $k$, we have with $\ell_k = |w_{[k,n]}| + 2$ that $\rho(x_k,\breve{c}) \leqslant 2^{-\ell_k}$.
If $\breve{c} = 0^{\infty} = 1^{\infty}$, assuming $v_{[1,2]} \notin \{00,11\}$ implies that $w_{[k,n]} v_{[1,2]} \notin \mc G$.
If $\breve{c} \neq 0$, we choose $v_{[1,2]}$  not to be a neighbor of $v'$, and hence $w_{[k,n]} v_{[1,2]}$ is not a neighbor of $w_{[k,n]}v'$, again implying $w_{[k,n]}v_{[1,2]} \notin \mc G$. Thus, $\rho(y_k,\breve{c}) \geqslant 2^{-\ell_k} \geqslant \rho(x_k,\breve{c})$, and we conclude that
\begin{equation}
\label{EQ:psi-variation-0}
\psi^c(x_k) - \psi^c(y_k) \leqslant 0
\end{equation}
holds for all $k=rp + 1$ with $0 \leqslant r < m$.
If $p=1$, this holds for all possible $k$ and we are done. Otherwise, we proceed with $k$ of the form $k = rp+ t$ with $0 \leqslant r < m-2$ and $2\leqslant t \leqslant p$. 
In this case, the minimality of the period $p$ implies that $w_{[k,n]}$ does not start with $u$ and therefore $w_{[k,k+p-1]}$ is not a prefix of $\breve{c}$. By Lemma~\ref{LEM:periodic-hitting-depth}, this means that the hitting depth in $wv$ can be at most $k+2p-2$. 
This means that by Lemma~\ref{LEM:psi-on-hitting} we obtain
\[
\psi^c(x_k) - \psi^c(y_k) \leqslant 2^{k+2p -n},
\]
and combining this with \eqref{EQ:psi-variation-0}
gives
\[
\sum_{k=1}^{(m-2)p - 1} \psi^c(x_k) - \psi^c(y_k)
\leqslant \sum_{k=1}^{mp - 2p - 1} 2^{k+2p -mp} \leqslant 1.
\]
Finally, we consider $(m-2)p \leqslant k \leqslant mp$. Applying Lemma~\ref{LEM:closing-extension-II} to the word $w_{[(m-2)p,n]} v_{[1,2]}$ we can choose $v$ in such a manner that $w_{[k,n]} v \notin \mc G$ for all such $k$ and $|v|$ is bounded by $\ell_0 = \max \{ C ,2\log_2 \log_{3/2} (2p + 3) \} + 2$. 
For all such $k$, we use the general bound in \eqref{EQ:psi-variation-general} to obtain
\[
\sum_{k = (m-2)p}^{mp} \psi^c(x_k) - \psi^c(y_k)
\leqslant (2p + 1) 2^{\ell_0 + 1}.
\]
The sum of all contributions can be therefore bounded in terms of an expression in $p$ which is smaller than $\varepsilon n$ for large enough $n$.
With $n$ large enough, the bound $\ell_0$ for $|v|$ is smaller than $2 \log_2 \log_{3/2} n$ and by the properties established above, we see that $w_{[k,n]}v \notin \mc G$ for all $1\leqslant k \leqslant n$.
\end{proof}

\end{document}
